\par\noindent\textbf{(1)} 选取正方向定义 \(J=\sigma(E-S\,\partial_xT)\) 和 \(J_Q=\Pi J-\kappa\partial_xT\)。开路 \(J=0\) 给 \(E=S\partial_xT\)，故两端电势差 \(\Delta V=-\int E\,dx=-\int_{T_1}^{T_2}S(T)\,dT\)，符号需与测量端定义一致。
\par\noindent\textbf{(2)} Peltier系数定义为均匀温度下每单位电流携带的可逆热流 \(\Pi=J_Q/J\)（扣除Fourier项后）。在不同材料界面，电流穿界时可逆Peltier热率 \( \dot Q=I(\Pi_A-\Pi_B)\) ，而内部焦耳项 \(J^2/\sigma\ge0\) 属不可逆耗散，不能合并进 \(\Pi\)。
\par\noindent\textbf{(3)} 令热力学力为 \(X_e=E/T\)、\(X_T=-\partial_xT/T^2\)，线性Onsager矩阵满足 \(L_{eQ}=L_{Qe}\)（无磁场）。与本构对比可得第一Kelvin关系 \(\Pi=ST\)。这是可逆热电交叉效应的互易约束，不由简单能量守恒单独保证。
\par\noindent\textbf{(4)} 沿温度梯度载流的可逆体热含Thomson项，系数 \(\tau=T\,dS/dT\)，称第二Kelvin关系；若 \(S\) 为温度常数则 \(\tau=0\)，但只要 \(S\neq0\) 仍有Peltier热。高温相变或非局域电子输运时线性系数可能不适用。
